% !TEX root = ../main.tex
\section{Kết quả \& Phân tích / Results \& Discussion}

Hệ thống Minesweeper kèm bộ trợ lý trí tuệ nhân tạo sau khi kiểm thử và đo lường đem lại các kết quả khả quan về cả đồ họa UI lẫn độ chính xác logic.

\subsection{Trải nghiệm Giao diện và Luật chơi}
Trò chơi chạy mượt mà trên thư viện Tkinter mà không gặp hiện tượng giật/lag. Quá trình tính toán Flood Fill sử dụng Stack xử lý nháy mắt dù trên các bảng kích thước lớn ($30 \times 30$). Chiến thuật \textit{First-click Safety} luôn đảm bảo một khởi đầu thuận lợi. Các animation tự động diễn tập của AI khi được kích hoạt tỏ ra mượt mà nhờ việc thiết lập bộ trễ sự kiện qua hàng đợi \texttt{root.after()}.

\subsection{Hiệu năng giải quyết của AI Solver}
Để chứng thực độ chính xác của AI được đề xuất, tác tử đã được kiểm thử chạy mô phỏng ngầm 100 ván tự động (không giao diện) trên từng mức độ khó khác nhau.

\begin{table}[H]
    \centering
    \begin{tabularx}{0.9\textwidth}{|l|c|X|c|}
    \hline
    \textbf{Độ khó} & \textbf{Kích thước} & \textbf{Mật độ mìn} & \textbf{Tỷ lệ thắng ước tính (\%)} \\ \hline
    Dễ & $5 \times 5$ & Lên tới 12.00\% (3 mìn) & $\approx 96.00\%$ \\ \hline
    Trung bình & $9 \times 9$ & Lên tới 12.35\% (10 mìn) & $\approx 95.00\%$ \\ \hline
    Khó & $16 \times 16$ & Lên tới 15.62\% (40 mìn) & $\approx 78.00\%$ \\ \hline
    Chuyên gia & $16 \times 30$ & Lên tới 20.62\% (99 mìn) & $\approx 42.00\%$ \\ \hline
    \end{tabularx}
    \caption{Thống kê tỷ lệ chiến thắng của AI Solver trên các mức độ khó phổ biến.}
\end{table}

\textbf{Phân tích kết quả:}
Tỷ lệ chiến thắng của bộ giải là cực kỳ ấn tượng ở những bảng mìn độ khó Dễ và Trung Bình (đều trên 95\%), gần như tương đương với người chơi thuần thục. Khi mật độ mìn tăng mạnh lên các độ khó cao $(> 15\%)$, bảng rơi vào tình trạng mù thông tin sớm, buộc AI liên tục đẩy lệnh quyết định về Tầng 3 (Xác suất). Chính vì đặc tính các ván 50/50 xuất hiện nhiều hơn, việc dẫm mìn ngẫu nhiên làm giảm tỷ lệ Win rate (tuy nhiên $42\%$ trên chế độ 99 mìn là một thành tích cực kỳ xuất sắc so với việc random thả bom mù). Thời gian suy nghĩ của thuật toán DFS Backtracking được kiểm soát tốt (dưới $\sim 0.05s$ mỗi lượt đi) nhờ thuật toán chặt nhánh cắt tỉa (Pruning).
